% ===========================================================================
%  Preamble for "The Engine's Mind" -- shared by the book and any excerpt.
% ===========================================================================
\usepackage[utf8]{inputenc}
\usepackage[T1]{fontenc}
\usepackage{lmodern}
\usepackage{microtype}
\usepackage[a4paper,margin=1in,headheight=14pt]{geometry}
\usepackage{amsmath,amssymb,amsthm}
\usepackage{booktabs}
\usepackage{array}
\usepackage{enumitem}
\usepackage{etoolbox}
\usepackage{fancyhdr}
\usepackage{xcolor}
\usepackage[most]{tcolorbox}
\usepackage{tikz}
\usetikzlibrary{positioning,calc,arrows.meta,shapes.geometric,fit,backgrounds}
\usepackage{xskak}
\usepackage{caption}
\usepackage[hidelinks,colorlinks=true,linkcolor=NavyBlue,urlcolor=NavyBlue,
            citecolor=NavyBlue]{hyperref}

% --------------------------------------------------------------------------
% Colours
% --------------------------------------------------------------------------
\definecolor{NavyBlue}{HTML}{1F3B73}
\definecolor{BuildBlue}{HTML}{2C5AA0}
\definecolor{IdeaGreen}{HTML}{2E6E4E}
\definecolor{PitfallRed}{HTML}{9B2C2C}
\definecolor{WorkGrey}{HTML}{4A4A4A}
\definecolor{DataTeal}{HTML}{1C6E72}
\definecolor{HypPurple}{HTML}{5B3A87}
\definecolor{PaleGrey}{HTML}{F4F4F6}

% --------------------------------------------------------------------------
% Headers
% --------------------------------------------------------------------------
\pagestyle{fancy}
\fancyhf{}
\fancyhead[LE]{\small\nouppercase{\leftmark}}
\fancyhead[RO]{\small\nouppercase{\rightmark}}
\fancyfoot[C]{\thepage}
\renewcommand{\headrulewidth}{0.3pt}

% --------------------------------------------------------------------------
% Math shorthands
% --------------------------------------------------------------------------
\DeclareMathOperator*{\argmax}{arg\,max}
\DeclareMathOperator*{\argmin}{arg\,min}
\DeclareMathOperator{\softmax}{softmax}
\DeclareMathOperator{\relu}{ReLU}
\DeclareMathOperator{\LN}{LayerNorm}
\DeclareMathOperator{\Ent}{H}
\newcommand{\E}{\mathbb{E}}
\newcommand{\R}{\mathbb{R}}
\newcommand{\given}{\,|\,}
\newcommand{\Qsa}{Q(s,a)}
\newcommand{\Nsa}{N(s,a)}
\newcommand{\Usa}{U(s,a)}
\newcommand{\Psa}{P(a \given s)}
\newcommand{\vv}[1]{\mathbf{#1}}

% --------------------------------------------------------------------------
% Boxes
% --------------------------------------------------------------------------
\newtcolorbox{keyidea}[1][]{
  breakable, enhanced, colback=IdeaGreen!5, colframe=IdeaGreen,
  boxrule=0.6pt, left=6pt, right=6pt, top=5pt, bottom=5pt,
  fonttitle=\bfseries, title={Key idea}, #1}

\newtcolorbox{pitfall}[1][]{
  breakable, enhanced, colback=PitfallRed!4, colframe=PitfallRed,
  boxrule=0.6pt, left=6pt, right=6pt, top=5pt, bottom=5pt,
  fonttitle=\bfseries, title={Where people go wrong}, #1}

\newtcolorbox{notationbox}[1][]{
  breakable, enhanced, colback=PaleGrey, colframe=WorkGrey!60,
  boxrule=0.5pt, left=6pt, right=6pt, fonttitle=\bfseries,
  title={Notation}, #1}

% Confidence tiers -- the honesty apparatus (see chapter 0).
\newtcolorbox{tier}[2][]{
  breakable, enhanced, colback=#2!4, colframe=#2, boxrule=0.6pt,
  left=6pt, right=6pt, top=4pt, bottom=4pt, fonttitle=\bfseries, #1}

\newenvironment{established}[1][]
  {\begin{tier}[title={Confidence: established}#1]{BuildBlue}}{\end{tier}}
\newenvironment{researchfinding}[1][]
  {\begin{tier}[title={Confidence: published research finding}#1]{DataTeal}}{\end{tier}}
\newenvironment{hypothesis}[1][]
  {\begin{tier}[title={Confidence: this project's hypothesis --- not yet verified}#1]{HypPurple}}{\end{tier}}
\newenvironment{realdata}[1][]
  {\begin{tier}[title={Real engine output}#1]{DataTeal}}{\end{tier}}

% --------------------------------------------------------------------------
% Formula-building steps: the spine of the book's method
% --------------------------------------------------------------------------
\newcounter{bstep}[chapter]
\newtcolorbox{buildbox}[1]{
  breakable, enhanced, colback=BuildBlue!4, colframe=BuildBlue,
  boxrule=0.6pt, left=6pt, right=6pt, top=4pt, bottom=4pt,
  fonttitle=\bfseries, title={#1}}
\newenvironment{attempt}[2]{%
  \begin{buildbox}{Attempt #1: #2}}{\end{buildbox}}
\newenvironment{finalform}[1][The formula, final form]{%
  \begin{buildbox}{#1}}{\end{buildbox}}

% --------------------------------------------------------------------------
% Worked examples and exercises
% --------------------------------------------------------------------------
\newcounter{wex}[chapter]
\renewcommand{\thewex}{\thechapter.\arabic{wex}}
\newenvironment{worked}[1]{%
  \refstepcounter{wex}%
  \par\medskip
  \begin{tcolorbox}[breakable, enhanced, colback=white, colframe=WorkGrey!70,
     boxrule=0.5pt, left=6pt, right=6pt, top=4pt, bottom=4pt,
     fonttitle=\bfseries, title={Worked example \thewex\ --- #1}]}%
  {\end{tcolorbox}\par\medskip}

\newcounter{exer}[chapter]
\renewcommand{\theexer}{\thechapter.\arabic{exer}}
% \exhead[short tag]{label}  -- label is used by \solutionto in appendix B
\newcommand{\exhead}[2][]{%
  \refstepcounter{exer}\label{ex:#2}\par\medskip\noindent
  \textbf{Exercise \theexer}\ifblank{#1}{}{\ \textit{(#1)}}\quad\ignorespaces}
\newenvironment{exercises}{%
  \section*{Exercises}\addcontentsline{toc}{section}{Exercises}}{}

\newcommand{\solutionto}[1]{%
  \par\medskip\noindent\textbf{Solution \ref{#1}.}\quad\ignorespaces}

% Chapter-opening "what you will be able to do" block
\newtcolorbox{objectives}[1][]{
  breakable, enhanced, colback=PaleGrey, colframe=NavyBlue,
  boxrule=0.7pt, left=8pt, right=8pt, top=6pt, bottom=6pt,
  fonttitle=\bfseries, title={After this chapter you can}, #1}

% Code / repo pointer
\newtcolorbox{repobox}[1][]{
  breakable, enhanced, colback=white, colframe=WorkGrey!50, boxrule=0.5pt,
  left=6pt, right=6pt, fonttitle=\bfseries\ttfamily,
  title={in this repository}, #1}

% --------------------------------------------------------------------------
% Chess board defaults
% --------------------------------------------------------------------------
\setchessboard{
  showmover=true,
  boardfontsize=15pt,
  labelfontsize=7pt,
  moverstyle=triangle,
  linewidth=0.4pt}
\newcommand{\minidiagram}[1]{\chessboard[setfen=#1,boardfontsize=13pt,labelfontsize=6pt]}

% Tighter lists
\setlist[itemize]{itemsep=2pt,topsep=3pt}
\setlist[enumerate]{itemsep=2pt,topsep=3pt}
